\documentclass[11pt]{article}

\usepackage[margin=1in]{geometry}
\usepackage[T1]{fontenc}
\usepackage{lmodern}

\usepackage{amsmath, amssymb, mathtools}

\usepackage{setspace}

\usepackage[most]{tcolorbox}
\usepackage{xcolor}

\tcbset{
  enhanced,
  boxrule=0pt,
  frame hidden,
  breakable
}

\usepackage{fancyhdr}
\pagestyle{fancy}
\fancyhf{}
\setlength{\headheight}{52pt}
\renewcommand{\headrulewidth}{0pt}

\newcommand{\Course}{Math 4610 Spring 2026}
\newcommand{\Instructor}{Homework 7}
\newcommand{\AssignmentType}{Homework}
\newcommand{\AssignmentNumber}{7}

\newcommand{\Contributors}{
Marks, Stephen
}

\fancyhead[L]{\Course\\ \Instructor}
\fancyhead[R]{\Contributors}

\newcommand{\ProblemDivider}[1]{%
  \vspace{1.0em}
  \noindent\textbf{#1}\par
  \vspace{0.35em}
}

\definecolor{problemBack}{RGB}{235,242,250}
\definecolor{problemFrame}{RGB}{70,110,160}
\definecolor{exerciseGray}{gray}{0.96}
\definecolor{exerciseGrayFrame}{gray}{0.45}

\newtcolorbox{problemBox}{
  colback=exerciseGray,
  colframe=exerciseGrayFrame,
  arc=2mm,
  left=8pt,right=8pt,top=8pt,bottom=8pt
}

\newtcolorbox{solutionBox}{
  colback=white,
  colframe=white,
  arc=0mm,
  left=0pt,right=0pt,top=0pt,bottom=0pt
}

\newcommand{\ProblemAnnounce}{\noindent}
\newcommand{\SolutionAnnounce}{\noindent\textbf{Solution.}\par\medskip}
\newcommand{\Exercise}[1]{\textbf{Exercise~#1.}\;}

\newcommand{\R}{\mathbb{R}}
\newcommand{\C}{\mathbb{C}}
\newcommand{\F}{\mathbb{F}}
\newcommand{\LL}{\mathcal{L}}
\newcommand{\spann}{\operatorname{span}}
\newcommand{\conjugate}{\operatorname{conjugate}}

\begin{document}

\begin{center}
  {\Large \textbf{\AssignmentType~\AssignmentNumber}}
\end{center}
\vspace{0.75em}

\ProblemDivider{1}

\begin{problemBox}
\ProblemAnnounce
\Exercise{7.A.1}
Suppose $n$ is a positive integer. Define $T \in \LL(\F^n)$ by
\[
T(z_1, \dots, z_n) = (0, z_1, \dots, z_{n-1}).
\]
Find a formula for $T^*(z_1, \dots, z_n)$.
\end{problemBox}

\begin{solutionBox}
\SolutionAnnounce
\doublespacing

Suppose $v,w \in \F^n$. Then
\[
\langle Tv,w\rangle = 0w_1 + v_1w_2 + v_2w_3 + \cdots + v_{n-1}w_n
\]
\[
= v_1w_2 + v_2w_3 + \cdots + v_{n-1}w_n
\]
\[
= v_1w_2 + v_2w_3 + \cdots + v_{n-1}w_n + 0v_n = \langle v,T^*w\rangle.
\]
Thus, $T^*w = (w_2,w_3,\dots,w_n,0)$

\end{solutionBox}

\ProblemDivider{2}

\begin{problemBox}
\ProblemAnnounce
\Exercise{7.A.3}
Suppose $T \in \LL(V)$ and $\lambda \in \F$. Prove that
\[
\lambda \text{ is an eigenvalue of } T \Longleftrightarrow \conjugate(\lambda) \text{ is an eigenvalue of } T^*.
\]
\end{problemBox}

\begin{solutionBox}
\SolutionAnnounce
\doublespacing

Suppose $\lambda$ is an eigenvalue of $T$ and $v \in V$. Then
\[
Tv = \lambda v
\]
\[
\Longleftrightarrow (T-\lambda I)v = 0.
\]
\[
\Longleftrightarrow (T-\lambda I) \text{ is not invertible}.
\]
\[
\Longleftrightarrow (T-\lambda I)^* \text{ is not invertible}.
\]
\[
\Longleftrightarrow (T^*-\conjugate(\lambda)I)w = 0 \text{ for some } w \neq 0 \in V.
\]
\[
\Longleftrightarrow T^*w = \conjugate(\lambda)w
\]

\end{solutionBox}

\ProblemDivider{3}

\begin{problemBox}
\ProblemAnnounce
\Exercise{7.A.4}
Suppose $T \in \LL(V)$ and $U$ is a subspace of $V$. Prove that
\[
U \text{ is invariant under } T \Longleftrightarrow U^\perp \text{ is invariant under } T^*.
\]
\end{problemBox}

\begin{solutionBox}
\SolutionAnnounce
\doublespacing

Suppose $U$ is invariant under $T$.
\[
U \text{ is invariant under } T.
\]
\[
\Longleftrightarrow Tu \in U \quad \forall u \in U.
\]
\[
\Longleftrightarrow Tu \in (U^\perp)^\perp = U \quad \forall u \in U.
\]
\[
\Longleftrightarrow \langle Tu, w\rangle = 0 \quad \forall u \in U, \forall w \in U^\perp.
\]
\[
\Longleftrightarrow \langle u, T^*w\rangle = 0 \quad \forall u \in U, \forall w \in U^\perp.
\]
\[
\Longleftrightarrow T^*w \in U^\perp \quad \forall w \in U^\perp.
\]
\[
\Longleftrightarrow U^\perp \text{ is invariant under } T^*.
\]

\end{solutionBox}

\ProblemDivider{4}

\begin{problemBox}
\ProblemAnnounce
\Exercise{7.B.1}
Prove that a normal operator on a complex inner product space is self-adjoint if and only if all its eigenvalues are real. This exercise strengthens the analogy (for normal operators) between self-adjoint operators and real numbers.
\end{problemBox}

\begin{solutionBox}
\SolutionAnnounce
\doublespacing

$(\rightarrow)$ Suppose $V$ is a complex inner product space and $T \in \LL(V)$ is normal. If $T$ is self-adjoint then every eigenvalue of $T$ is real by 7.12.

$(\leftarrow)$ Suppose that every eigenvalue of $T$ is real. 7.31 implies that there is an orthonormal basis of $V$ with respect to which the matrix $\mathcal{M}(T)$ is diagonal. Because each eigenvalue of $T$ is real, the diagonal entries of $\mathcal{M}(T)$ must be real. It follows that $\mathcal{M}(T)$ equals its own conjugate transpose and thus that $T$ is self-adjoint.

\end{solutionBox}

\ProblemDivider{5}

\begin{problemBox}
\ProblemAnnounce
\Exercise{7.B.2}
Suppose $\F = \C$. Suppose $T \in \LL(V)$ is normal and has only one eigenvalue. Prove that $T$ is a scalar multiple of the identity operator.
\end{problemBox}

\begin{solutionBox}
\SolutionAnnounce
\doublespacing

Suppose that $\lambda \in \C$ is the only eigenvalue of $T$. Then there is an orthonormal basis of $V$ with respect to which $\mathcal{M}(T)$ is diagonal by 7.31. Since $\lambda$ is the only eigenvalue of $T$, each diagonal entry of $\mathcal{M}(T)$ must be equal to $\lambda$. Thus $T = \lambda I$.

\end{solutionBox}

\ProblemDivider{6}

\begin{problemBox}
\ProblemAnnounce
\Exercise{7.B.3}
Suppose $\F = \C$ and $T \in \LL(V)$ is normal. Prove that the set of eigenvalues of $T$ is contained in $\{0, 1\}$ if and only if there is a subspace $U$ of $V$ such that $T = P_U$.
\end{problemBox}

\begin{solutionBox}
\SolutionAnnounce
\doublespacing

$(\leftarrow)$ If there exists a subspace $U$ of $V$ such that $T = P_U$ then $T^2 = T$ then it follows from 5.A.8 that the set of eigenvalues of $T$ is contained in $\{0, 1\}$.

$(\rightarrow)$ Suppose that the set of eigenvalues of $T$ is contained in $\{0, 1\}$. Then there is an orthonormal basis $\mathcal{B}$ of $V$ consisting of eigenvectors of $T$ by 7.31. Each basis vector in $\mathcal{B}$ must correspond either to the eigenvalue $0$ or the eigenvalue $1$. Let $u_1, \dots, u_m$ be those basis vectors in $\mathcal{B}$ corresponding to the eigenvalue $1$ and let $v_1, \dots, v_n$ be those basis vectors in $\mathcal{B}$ corresponding to the eigenvalue $0$, but either of these lists may be empty. Let
\[
U = \spann(u_1, \dots, u_m).
\]
Because $\mathcal{B}$ is an orthonormal basis of $V$, it follows that
\[
U^\perp = \spann(v_1, \dots, v_n)
\]
and thus
\[
P_Uu_k = u_k = Tu_k
\]
and
\[
P_Uv_k = 0 = Tv_k.
\]
Thus $T = P_U$.

\end{solutionBox}

\end{document}
